% 论文骨架：无官方模板且中文写作时的默认起点（xelatex + ctex）。
% 有官方模板时以官方 class 为准，不要用本文件覆盖官方模板。
% 编译：latexmk -xelatex -halt-on-error -interaction=nonstopmode main.tex
%       无 latexmk 时：xelatex main.tex ×2（有 bib 则 xelatex → biber → xelatex ×2）
\documentclass[12pt,a4paper]{ctexart}

\usepackage[margin=2.5cm]{geometry}
\usepackage{amsmath,amssymb}
\usepackage{graphicx}
\usepackage{booktabs}
\usepackage{tabularx}
\usepackage{caption}
\usepackage{listings}
\usepackage[hidelinks]{hyperref}

\graphicspath{{../图/}}
\captionsetup{font=small}

% 数字宏：由 scripts/ledger.py --emit-tex 从 results_ledger.json 生成。
% 正文只引用宏，禁止手抄数字；本文件不存在时先跑一次生成命令。
\IfFileExists{numbers.tex}{\input{numbers.tex}}{\newcommand{\NUMBERSMISSING}{}}

\title{题号 + 题目名称}
\author{}                       % 匿名：不写姓名与学校
\date{}

\begin{document}
\maketitle
\thispagestyle{empty}

\begin{abstract}
逐问一句「方法 + 关键数值结论」，数字一律引用宏（如 \verb|\qOneMainMetric|），
机会约束类结论必须同时给出置信下界。最后一句总括模型特色。
摘要是评委唯一必读页；章程要求一页时用 \texttt{scripts/latex\_gate.py --aux} 程序化核实。

\textbf{关键词：} 关键词一；关键词二；关键词三
\end{abstract}

\newpage
\setcounter{page}{1}

\section{问题重述}
简明复述，不抄题面全文。

\section{模型假设}
每条假设附合理性依据；口径决策表的披露版放这里（采用口径 + 被否口径的定量归谬一句话）。

\section{符号说明}
\begin{table}[htbp]
  \centering
  \caption{符号说明}
  \begin{tabularx}{0.9\textwidth}{llX}
    \toprule
    符号 & 单位 & 含义 \\
    \midrule
    $x$ & -- & 示例 \\
    \bottomrule
  \end{tabularx}
\end{table}

\section{问题一：分析与建模}
\subsection{问题分析与准备}
\subsection{模型建立与求解}
\subsection{结果检验}
证书、收敛诊断、与独立核验路线的互证；差异 $\gtrsim \Delta$ 必须解释。

\begin{figure}[htbp]
  \centering
  % \includegraphics[width=0.8\linewidth]{q1_main.pdf}
  \caption{图题自足：不看正文也能说明这张图证明了什么（非等比缩放须标注比例）。}
  \label{fig:q1main}
\end{figure}

\section{问题二：分析与建模}

\section{灵敏度与稳健性分析}

\section{模型优缺点与改进方向}
缺点写实：实际达到的证据等级（S1--S4）、残余间隙、未排除的区域。

\begin{thebibliography}{9}
  \bibitem{ref1} 作者. 文献题名[J]. 期刊, 年份, 卷(期): 页码.
\end{thebibliography}

\appendix
\section{附录：明细表与推导过程}
正文只放证据链主线，明细表、逐点扫描、推导全过程放这里——这是在不砍证据的前提下压正文页数的唯一手段。

\section{附录：核心代码与复现说明}
\begin{lstlisting}[basicstyle=\ttfamily\small,breaklines=true]
# 复现命令示例
python run_all.py --stage smoke
\end{lstlisting}

\section{附录：AI 工具使用声明}
按赛事章程规定的格式与位置如实声明（漏报可取消评奖资格）。

\end{document}
